\ifdefined\gkver\else\def\gkver{student}\fi
\documentclass[\gkver]{gaokaozhenti}
\begin{document}

\chapter{2022年辽宁}

\begin{enumerate}
\item
%% number: 1
%% paperName: 2022年辽宁高考第 1 题 4 分
%% typeId: 1
%% type: 单选题
%% chapter: 直线运动
%% point: 路程与位移
%% method: 无
%% score: 4
%% degree: 950
%% duplicateId: 0
%% body:
如图所示，桥式起重机主要由可移动“桥架”“小车”和固定“轨道”三部分组成。在某次作业中桥架沿轨道单向移动了 $8 \Um$，小车在桥架上单向移动了 $6 \Um$。该次作业中小车相对地面的位移大小为\xzanswer{C}

\onepicture[w=6.80cm]{figs/01.png}

\fourchoices[answer=C]
{$6 \Um$}
{$8 \Um$}
{$10 \Um$}
{$14 \Um$}

%% answer: C
%% memo:
\memoanswer{
根据位移是起点到终点的有向线段，可知，该次作业中小车相对地面的位移为两个相互垂直方向移动的分位移矢量和，即该次作业中小车相对地面的位移大小为 $x = \sqrt{x_1^2 + x_2^2} = \sqrt{8^2 + 6^2} \Um = 10 \Um$，选项 C 正确。
}

\item
%% number: 2
%% paperName: 2022年辽宁高考第 2 题 4 分
%% typeId: 1
%% type: 单选题
%% chapter: 原子和原子核
%% point: 质能方程
%% method: 无
%% score: 4
%% degree: 850
%% duplicateId: 0
%% body:
2022年1月，中国锦屏深地实验室发表了首个核天体物理研究实验成果。表明我国核天体物理研究已经跻身国际先进行列。实验中所用核反应方程为 \ce{X + ^23_12Mg -> ^26_13Al}，已知 X、\ce{^23_12Mg}、\ce{^26_13Al} 的质量分别为 $m_{1}$、$m_{2}$、$m_{3}$，真空中的光速为 $c$，该反应中释放的能量为 $E$。下列说法正确的是\xzanswer{B}

\fourchoices[answer=B]
{X 为氘核 \ce{^2_1H}}
{X 为氚核 \ce{^3_1H}}
{$E = (m_{1} + m_{2} + m_{3})c^{2}$}
{$E = (m_{1} + m_{2} - m_{3})c^{2}$}

%% answer: B
%% memo:
\memoanswer{
根据核反应遵循的质量数守恒和电荷数守恒可知，X 的质量数为 3，电荷数为 1，为氚核 \ce{^3_1H}，选项 A 错误，B 正确；

因该反应为人工转变，反应前两种粒子都有动能（总动能设为 $E_{\mathrm{k}1}$），反应后的生成物也有动能 $E_{\mathrm{k}2}$，根据质能方程可知，由于质量亏损反应放出的能量为 $\Delta E = \Delta m c^{2} = (m_{1} + m_{2} - m_{3}) c^{2}$，则反应释放的能量为

$E = E_{\mathrm{k}1} + \Delta E - E_{\mathrm{k}2} = E_{\mathrm{k}1} - E_{\mathrm{k}2} + \left(m_{1} + m_{2} - m_{3}\right)c^{2}$，选项 CD 错误。
}

\item
%% number: 3
%% paperName: 2022年辽宁高考第 3 题 4 分
%% typeId: 1
%% type: 单选题
%% chapter: 机械波
%% point: 波的图象
%% method: 无
%% score: 4
%% degree: 750
%% duplicateId: 0
%% body:
一列简谐横波沿 $x$ 轴正方向传播，某时刻的波形如图所示，关于质点 P 的说法正确的是\xzanswer{A}

\onepicture[w=4.40cm]{figs/03.png}

\fourchoices[answer=A]
{该时刻速度沿 $y$ 轴正方向}
{该时刻加速度沿 $y$ 轴正方向}
{此后 $\frac{1}{4}$ 周期内通过的路程为 $A$}
{此后 $\frac{1}{2}$ 周期内沿 $x$ 轴正方向迁移为 $\frac{1}{2}\lambda$}

%% answer: A
%% memo:
\memoanswer{
波沿 $x$ 轴正向传播，由“同侧法”可知，该时刻质点 P 的速度方向沿 $y$ 轴正向，加速度沿 $y$ 轴负向，选项 A 正确，B 错误。

在该时刻质点 P 不在特殊位置，则在 $\frac{1}{4}$ 周期内的路程不一定等于 $A$，选项 C 错误；

质点只能在自己平衡位置附近振动，而不随波迁移，选项 D 错误。
}

\item
%% number: 4
%% paperName: 2022年辽宁高考第 4 题 4 分
%% typeId: 1
%% type: 单选题
%% chapter: 力物体的平衡
%% point: 三力平衡
%% method: 无
%% score: 4
%% degree: 650
%% duplicateId: 0
%% body:
如图所示，蜘蛛用蛛丝将其自身悬挂在水管上，并处于静止状态。蛛丝 OM、ON 与竖直方向夹角分别为 $\alpha$、$\beta$（$\alpha > \beta$）。用 $F_{1}$、$F_{2}$ 分别表示 OM、ON 的拉力，则\xzanswer{D}

\onepicture[w=4.40cm]{figs/04.png}

\fourchoices[answer=D]
{$F_{1}$ 的竖直分力大于 $F_{2}$ 的竖直分力}
{$F_{1}$ 的竖直分力等于 $F_{2}$ 的竖直分力}
{$F_{1}$ 的水平分力大于 $F_{2}$ 的水平分力}
{$F_{1}$ 的水平分力等于 $F_{2}$ 的水平分力}

%% answer: D
%% memo:
\memoanswer{
此题的模型属于死结，对结点 O 受力分析可得，水平方向 $F_{1}\sin\alpha = F_{2}\sin\beta$

即 $F_{1}$ 的水平分力等于 $F_{2}$ 的水平分力，选项 C 错误，D 正确；

对结点 O 受力分析可得，竖直方向 $F_{1}\cos\alpha + F_{2}\cos\beta = mg$，解得 $F_{1} = \frac{mg\sin\beta}{\sin(\alpha + \beta)}$，$F_{2} = \frac{mg\sin\alpha}{\sin(\alpha + \beta)}$

则 $F_{1}$ 的竖直分量 $F_{1y} = \frac{mg \sin \beta \cos \alpha}{\sin (\alpha + \beta)}$，$F_{2}$ 的竖直分量 $F_{2y} = \frac{mg \sin \alpha \cos \beta}{\sin (\alpha + \beta)}$。$F_{2}$ 的竖直分力 $F_{2y}$ 与 $F_{1}$ 的竖直分力 $F_{1y}$ 之差

$F_{2y} - F_{1y} = \frac{mg \sin \alpha \cos \beta}{\sin (\alpha + \beta)} - \frac{mg \sin \beta \cos \alpha}{\sin (\alpha + \beta)} = \frac{mg (\sin \alpha \cos \beta - \sin \beta \cos \alpha)}{\sin (\alpha + \beta)}$

因 $\sin \alpha \cos \beta - \cos \alpha \sin \beta = \sin (\alpha - \beta) > 0$

可知 $F_{2y} > F_{1y}$，选项 AB 错误。
}

\item
%% number: 5
%% paperName: 2022年辽宁高考第 5 题 4 分
%% typeId: 1
%% type: 单选题
%% chapter: 光学
%% point: 全反射
%% method: 无
%% score: 4
%% degree: 650
%% duplicateId: 0
%% body:
完全失重时，液滴呈球形，气泡在液体中将不会上浮。2021 年 12 月，在中国空间站“天宫课堂”的水球光学实验中，航天员向水球中注入空气形成了一个内含气泡的水球。如图所示，若气泡与水球同心，在过球心 O 的平面内，用单色平行光照射这一水球。下列说法正确的是\xzanswer{C}

\onepicture[w=4.60cm]{\includesvg{figs/05}}

\fourchoices[answer=C]
{此单色光从空气进入水球，频率一定变大}
{此单色光从空气进入水球，频率一定变小}
{若光线 1 在 M 处发生全反射，光线 2 在 N 处一定发生全反射}
{若光线 2 在 N 处发生全反射，光线 1 在 M 处一定发生全反射}

%% answer: C
%% memo:
\memoanswer{
由于光的频率是由光源决定的，与介质无关，所以此单色光从空气进入水球，频率不变，选项 AB 错误；

由光路图可看出光线 1 入射到水球的入射角小于光线 2 入射到水球的入射角，则光线 1 在水球外表面折射后的折射角小于光线 2 在水球外表面折射后的折射角，设水球半径为 $R$、气泡半径为 $r$、光线经过水球后的折射角为 $\alpha$、光线进入气泡的入射角为 $\theta$，根据几何关系，由正弦定理可得 $\frac{\sin (\pi - \theta)}{R} = \frac{\sin \alpha}{r}$

由此可知光线 2 的 $\theta$ 大于光线 1 的 $\theta$，故若光线 1 在 M 处发生全反射，光线 2 在 N 处一定发生全反射，C 正确、D 错误。
}

\item
%% number: 6
%% paperName: 2022年辽宁高考第 6 题 4 分
%% typeId: 1
%% type: 单选题
%% chapter: 气体性质
%% point: 热力学第一定律
%% method: 无
%% score: 4
%% degree: 650
%% duplicateId: 0
%% body:
一定质量的理想气体从状态 a 变化到状态 b，其体积 $V$ 和热力学温度 $T$ 变化图像如图所示，此过程中该系统\xzanswer{A}

\onepicture[w=4.40cm]{figs/06.png}

\fourchoices[answer=A]
{对外界做正功}
{压强保持不变}
{向外界放热}
{内能减少}

%% answer: A
%% memo:
\memoanswer{
理想气体从状态 a 变化到状态 b，体积增大，理想气体对外界做正功，选项 A 正确；

由题图 $V-T$ 图像可知 $V = V_{0} + kT$，根据理想气体的状态方程有 $\frac{pV}{T} = C$，联立有 $p = \frac{C}{k + \frac{V_{0}}{T}}$，可看出 $T$ 增大，$p$ 增大，选项 B 错误；

理想气体从状态 a 变化到状态 b，温度升高，内能增大，选项 D 错误；

理想气体从状态 a 变化到状态 b，由选项 AD 可知，理想气体对外界做正功且内能增大，则根据热力学第一定律可知气体向外界吸收热量，选项 C 错误。
}

\item
%% number: 7
%% paperName: 2022年辽宁高考第 7 题 4 分
%% typeId: 1
%% type: 单选题
%% chapter: 运动和力
%% point: 牛顿定律的应用
%% method: 无
%% score: 4
%% degree: 650
%% duplicateId: 0
%% body:
如图所示，一小物块从长 $1 \Um$ 的水平桌面一端以初速度 $v_{0}$ 沿中线滑向另一端，经过 $1 \Us$ 从另一端滑落。物块与桌面间动摩擦因数为 $\mu$，$g$ 取 $10 \Umsq$。下列 $v_{0}$、$\mu$ 值可能正确的是\xzanswer{B}

\onepicture[w=5.20cm]{\includesvg{figs/07}}

\fourchoices[answer=B]
{$v_{0} = 2.5 \Ums$}
{$v_{0} = 1.5 \Ums$}
{$\mu = 0.28$}
{$\mu = 0.25$}

%% answer: B
%% memo:
\memoanswer{
物块水平沿水平桌面中线做匀减速直线运动，则 $x = \frac{v_0 + v}{2}t$，由题述可知，$x = 1 \Um$，$t = 1 \Us$，$v > 0$，代入数据有 $v_{0} < 2 \Ums$，故 A 不可能，B 可能；

对物块做受力分析，由牛顿第二定律有 $a = -\mu g$，利用匀变速直线运动规律得 $v^{2} - v_{0}^{2} = 2ax$，整理有 $v_{0}^{2} - 2ax > 0$，由于 $v_{0} < 2 \Ums$ 可得 $\mu < 0.2$，故选项 CD 不可能。
}

\item
%% number: 8
%% paperName: 2022年辽宁高考第 8 题 4 分
%% typeId: 1
%% type: 单选题
%% chapter: 磁场
%% point: 带电粒子在磁场中的运动
%% method: 无
%% score: 4
%% degree: 650
%% duplicateId: 0
%% body:
粒子物理研究中使用的一种球状探测装置横截面的简化模型如图所示。内圆区域有垂直纸面向里的匀强磁场，外圆是探测器。两个粒子先后从 P 点沿径向射入磁场，粒子 1 沿直线通过磁场区域后打在探测器上的 M 点。粒子 2 经磁场偏转后打在探测器上的 N 点。装置内部为真空状态，忽略粒子重力及粒子间相互作用力。下列说法正确的是\xzanswer{AD}

\onepicture[w=4.80cm]{figs/08.png}

\fourchoices[answer=AD]
{粒子 1 可能为中子}
{粒子 2 可能为电子}
{若增大磁感应强度，粒子 1 可能打在探测器上的 Q 点}
{若增大粒子入射速度，粒子 2 可能打在探测器上的 Q 点}

%% answer: AD
%% memo:
\memoanswer{
由题图可看出粒子 1 没有偏转，说明粒子 1 不受磁场力作用，即粒子 1 不带电，则粒子 1 可能为中子；粒子 2 向上偏转，根据左手定则可知粒子 2 应该带正电，选项 A 正确、B 错误；

由以上分析可知粒子 1 为中子，则无论如何增大磁感应强度，粒子 1 都不会偏转，选项 C 错误；

粒子 2 在磁场中运动，洛伦兹力提供向心力有 $qvB = m\frac{v^{2}}{r}$，解得粒子运动轨迹半径 $r=\frac{mv}{qB}$。可知若增大粒子入射速度，则粒子 2 的半径增大，粒子 2 可能打在探测器上的 Q 点，选项 D 正确。
}

\item
%% number: 9
%% paperName: 2022年辽宁高考第 9 题 6 分
%% typeId: 2
%% type: 多选题
%% chapter: 曲线运动
%% point: 天体运动
%% method: 无
%% score: 6
%% degree: 550
%% duplicateId: 0
%% body:
如图所示，行星绕太阳的公转可以看成匀速圆周运动。在地图上容易测得地球—水星连线与地球—太阳连线夹角 $\alpha$，地球—金星连线与地球—太阳连线夹角 $\beta$，两角最大值分别为 $\alpha_{m}$、$\beta_{m}$，则\xzanswer{BC}

\onepicture[w=4.80cm]{figs/09.png}

\fourchoices[answer=BC]
{水星的公转周期比金星的大}
{水星的公转向心加速度比金星的大}
{水星与金星的公转轨道半径之比为 $\sin\alpha_{m}:\sin\beta_{m}$}
{水星与金星的公转线速度之比为 $\sqrt{\sin\alpha_{m}}:\sqrt{\sin\beta_{m}}$}

%% answer: BC
%% memo:
\memoanswer{
天体运动，万有引力提供向心力有，

$G\frac{Mm}{R^{2}}=m\frac{4\pi^{2}}{T^{2}}R=ma$

可解得 $T = 2\pi \sqrt{\frac{R^{3}}{GM}}$，$a = \frac{GM}{R^{2}}$。因为水星的公转半径比金星小，故可知水星的公转周期比金星小；水星的公转向心加速度比金星的大，故选项 A 错误，B 正确；

设水星的公转半径为 $R_{\text{水}}$，地球的公转半径为 $R_{\text{地}}$，当 $\alpha$ 角最大时有 $\sin\alpha_{m}=\frac{R_{\text{水}}}{R_{\text{地}}}$，同理可知有 $\sin\beta_{m}=\frac{R_{\text{金}}}{R_{\text{地}}}$

所以水星与金星的公转半径之比为 $R_{\text{水}}:R_{\text{金}}=\sin\alpha_{m}:\sin\beta_{m}$，选项 C 正确；

根据万有引力提供向心力，$G\frac{Mm}{R^{2}}=m\frac{v^{2}}{R}$，可得 $v=\sqrt{\frac{GM}{R}}$

由 $v = \sqrt{\frac{GM}{R}}$，$R_{\text{水}}:R_{\text{金}} = \sin\alpha_{m} : \sin\beta_{m}$ 联立解得 $v_{\text{水}}:v_{\text{金}} = \sqrt{\sin\beta_{m}} : \sqrt{\sin\alpha_{m}}$，选项 D 错误。
}

\item
%% number: 10
%% paperName: 2022年辽宁高考第 10 题 6 分
%% typeId: 2
%% type: 多选题
%% chapter: 电场
%% point: 带电粒子的直线运动
%% method: 无
%% score: 6
%% degree: 450
%% duplicateId: 0
%% body:
如图所示，带电荷量为 $6Q$（$Q > 0$）的球 1 固定在倾角为 $30{^\circ}$ 的光滑绝缘斜面上的 a 点，其正上方 $L$ 处固定一电荷量为 $-Q$ 的球 2，斜面上距 a 点 $L$ 处的 b 点有质量为 $m$ 的带电球 3，球 3 与一端固定的绝缘轻质弹簧相连并在 b 点处于静止状态。此时弹簧的压缩量为 $\frac{L}{2}$，球 2、3 间的静电力大小为 $\frac{mg}{2}$。迅速移走球 1 后，球 3 沿斜面向下运动。$g$ 为重力加速度，球的大小可忽略，下列关于球 3 的说法正确的是\xzanswer{BCD}

\onepicture[w=5.20cm]{figs/10.png}

\fourchoices[answer=BCD]
{带负电}
{运动至 a 点的速度大小为 $\sqrt{gL}$}
{运动至 a 点的加速度大小为 $2g$}
{运动至 ab 中点时对斜面的压力大小为 $\frac{3\sqrt{3}-4}{6}mg$}

%% answer: BCD
%% memo:
\memoanswer{
根据题图几何关系可知三小球构成一个边长为 $L$ 的等边三角形，小球 1 和 3 之间的力大于小球 2 和 3 之间的力，弹簧处于压缩状态，故小球 1 和 3 之间的作用力一定是斥力，

由于小球 1 带正电，所以小球 3 带正电，选项 A 错误；

小球 3 运动至 a 点时，弹簧的伸长量等于 $\frac{L}{2}$，根据对称性可知，小球 2 对小球 3 做功为 0；弹簧弹力做功为 0，故根据动能定理有 $mgL\sin\theta=\frac{1}{2}mv^{2}$，解得 $v=\sqrt{gL}$，选项 B 正确；

小球 3 在 b 点时，设小球 3 的电荷量为 $q$，根据库仑定律和平衡条件，得 $k\frac{qQ}{L^{2}}=\frac{mg}{2}$。

设弹簧的弹力为 $F$，根据受力平衡，沿斜面方向有 $F = k\frac{q \cdot 6Q}{L^{2}} - k\frac{qQ}{L^{2}}\sin 30^{\circ} - mg\sin 30^{\circ}$，

解得 $F=\frac{9}{4}mg$。小球 3 运动至 a 点时，弹簧的伸长量等于 $\frac{L}{2}$，根据对称性，利用牛顿第二定律可得 $F + k \frac{qQ}{L^{2}} \sin 30^{\circ} - mg \sin 30^{\circ} = ma$，解得 $a = 2g$，选项 C 正确；

当小球 3 运动至 ab 中点时，弹簧弹力为 0，此时小球 2 对小球 3 的力为

$F_{23}=k\frac{qQ}{(\frac{\sqrt{3}}{2}L)^{2}}=\frac{4}{3}\cdot k\frac{qQ}{L^{2}}=\frac{4}{3}\times\frac{mg}{2}=\frac{2}{3}mg$

斜面对小球的支持力为 $F_{N}=mg\cos30^{\circ}-F_{23}=\frac{\sqrt{3}}{2}mg-\frac{2}{3}mg=\frac{3\sqrt{3}-4}{6}mg$

根据牛顿第三定律可知，小球对斜面的压力大小为 $\frac{3\sqrt{3}-4}{6}mg$，选项 D 正确。
}

\item
%% number: 11
%% paperName: 2022年辽宁高考第 11 题 6 分
%% typeId: 4
%% type: 实验题
%% chapter: 电路
%% point: 电表改装
%% method: 无
%% score: 6
%% degree: 650
%% duplicateId: 0
%% body:
某同学要将一小量程电流表（满偏电流为 $250 \UuA$，内阻为 $1.2 \UkO$）改装成有两个量程的电流表，设计电路如图（\ref{2022辽宁11a}）所示，其中定值电阻 $R_{1} = 40 \UO$，$R_{2} = 360 \UO$。

\begin{enumerate}
	\item
	当开关 S 接 A 端时，该电流表的量程为 $0 \sim$ \tkanswer{1} $\UmA$；
	\item
	当开关 S 接 B 端时，该电流表的量程比接在 A 端时\tkanswer{大}（填“大”或“小”）；
	\item
	该同学选用量程合适的电压表（内阻未知）和此改装电流表测量未知电阻 $R_{x}$ 的阻值，设计了图（\ref{2022辽宁11b}）中两个电路。不考虑实验操作中的偶然误差，则使用\tkanswer{乙}（填“甲”或“乙”）电路可修正由电表内阻引起的实验误差。
\end{enumerate}

\twopicture[labelA=2022辽宁11a, labelB=2022辽宁11b, showlabel=false, wA=3.60cm, wB=8.60cm, align=right]
{figs/11a.png}
{figs/11b.png}

%% answer:
\jdanswer{
\begin{enumerate}
	\item
	1

	\item
	大

	\item
	乙
\end{enumerate}
}
%% memo:
\memoanswer{
（1）由图可知当 S 接 A 时，$R_{1}$ 和 $R_{2}$ 串联接入电路，和电流表并联，满偏时电流表两端的电压为 $U_{m}=I_{m}r=250\times10^{-6}\times1.2\times10^{3}\UV=0.3\UV$

此时 $R_{1}$ 和 $R_{2}$ 的电流为

$I = \frac{U_{\mathrm{m}}}{R_{1} + R_{2}} = \frac{0.3}{40 + 360}\UA = 0.75 \times 10^{-3}\UA = 0.75\UmA$

所以总电流为 $I_{\text{总}} = I_{m} + I = 1\UmA$，即量程为 $0 \sim 1\UmA$。

（2）当开关 S 接 B 端时，由图可知电阻 $R_{1}$ 和电流表串联再和电阻 $R_{2}$ 并联，由于和电流表并联的电阻变小，当电流表满偏时，流过电阻 $R_{2}$ 的电流变大，干路电流变大，即量程变大；所以开关 S 接 B 端时，该电流表的量程比接在 A 端时大。

（3）图甲是电流表的外接法，误差是由于电压表的分流引起的；图乙是电流表的内接法，误差是由于电流表的分压引起的，因为题目中电压表电阻未知，故采用图乙的方法可以修正由电表内阻引起的实验误差。
}

\item
%% number: 12
%% paperName: 2022年辽宁高考第 12 题 8 分
%% typeId: 4
%% type: 实验题
%% chapter: 直线运动
%% point: 自由落体运动
%% method: 无
%% score: 8
%% degree: 550
%% duplicateId: 0
%% body:
某同学利用如图所示的装置测量重力加速度，其中光栅板上交替排列着等宽度的遮光带和透光带（宽度用 $d$ 表示）。实验时将光栅板置于光电传感器上方某高度，令其自由下落穿过光电传感器。光电传感器所连接的计算机可连续记录遮光带、透光带通过光电传感器的时间间隔 $\Delta t$。

\begin{enumerate}
	\item
	除图中所用的实验器材外，该实验还需要\tkanswer{刻度尺}（填“天平”或“刻度尺”）；
	\item
	该同学测得遮光带（透光带）的宽度为 $4.50 \Ucm$，记录时间间隔的数据如表所示，

	\begin{table}[htbp]
	\centering
	\begin{tblr}{|c|c|c|c|c|}
	\hline
	编号 & 1遮光带 & 2遮光带 & 3遮光带 & $\ldots$ \\
	\hline
	$\Delta t/(\times 10^{-3}\Us)$ & 73.04 & 38.67 & 30.00 & $\ldots$ \\
	\hline
	\end{tblr}
	\end{table}

	根据上述实验数据，可得编号为 3 的遮光带通过光电传感器的平均速度大小为 $v_{3} =$ \tkanswer{1.5} $\Ums$（结果保留两位有效数字）；
	\item
	某相邻遮光带和透光带先后通过光电传感器的时间间隔为 $\Delta t_{1}$、$\Delta t_{2}$，则重力加速度 $g =$ \tkanswer{$\frac{2d(\Delta t_{1} - \Delta t_{2})}{\Delta t_{1}\Delta t_{2}(\Delta t_{1} + \Delta t_{2})}$}（用 $d$、$\Delta t_{1}$、$\Delta t_{2}$ 表示）；
	\item
	该同学发现所得实验结果小于当地的重力加速度，请写出一条可能的原因：\tkanswer{光栅板下落受到空气阻力的作用}。
\end{enumerate}

\onepicture[w=5.20cm, align=right]{figs/12.png}

%% answer:
\jdanswer{
\begin{enumerate}
	\item
	刻度尺

	\item
	1.5

	\item
	$\frac{2d(\Delta t_{1} - \Delta t_{2})}{\Delta t_{1}\Delta t_{2}(\Delta t_{1} + \Delta t_{2})}$

	\item
	光栅板下落受到空气阻力的作用
\end{enumerate}
}
%% memo:
\memoanswer{
（1）该实验测量重力加速度，不需要天平测质量；需要用刻度尺测量遮光带（透光带）的宽度，故需要刻度尺；

（2）根据平均速度的计算公式可知，编号为 3 的遮光带通过光电传感器的平均速度大小为 $v_{3} = \frac{d}{\Delta t} = \frac{4.5 \times 10^{-2} \Um}{30 \times 10^{-3} \Us} = 1.5 \Ums$。

（3）由于 $d$ 很小，可以用遮光带通过光电传感器的平均速度代替瞬时速度，所以有

$v_{1}=\frac{d}{\Delta t_{1}}$，$v_{2}=\frac{d}{\Delta t_{2}}$

$v_{2}=v_{1}+g\left(\frac{\Delta t_{2}+\Delta t_{1}}{2}\right)$

由匀变速直线运动规律，可得 $g = \frac{2d(\Delta t_{1} - \Delta t_{2})}{\Delta t_{1}\Delta t_{2}(\Delta t_{1} + \Delta t_{2})}$。

（4）由于光栅板下落过程中受到空气阻力的影响，所以竖直向下的加速度小于重力加速度。
}

\item
%% number: 13
%% paperName: 2022年辽宁高考第 13 题 10 分
%% typeId: 6
%% type: 计算题
%% chapter: 曲线运动
%% point: 圆周运动的应用
%% method: 无
%% score: 10
%% degree: 650
%% duplicateId: 0
%% body:
2022年北京冬奥会短道速滑混合团体 2 000 米接力决赛中，我国短道速滑队夺得中国队在本届冬奥会的首金。

\begin{enumerate}
	\item
	如果把运动员起跑后进入弯道前的过程看作初速度为零的匀加速直线运动，若运动员加速到速度 $v = 9 \Ums$ 时，滑过的距离 $x = 15 \Um$，求加速度的大小；
	\item
	如果把运动员在弯道滑行的过程看作轨道为半圆的匀速圆周运动，如图所示，若甲、乙两名运动员同时进入弯道，滑行半径分别为 $R_{\text{甲}} = 8 \Um$、$R_{\text{乙}} = 9 \Um$，滑行速率分别为 $v_{\text{甲}} = 10 \Ums$、$v_{\text{乙}} = 11 \Ums$，求甲、乙过弯道时的向心加速度大小之比，并通过计算判断哪位运动员先出弯道。
\end{enumerate}

\onepicture[w=3.20cm, align=right]{figs/13.png}

%% answer:
\jdanswer{
\begin{enumerate}
	\item
	$2.7 \Umsq$

	\item
	$\frac{225}{242}$，甲
\end{enumerate}
}
%% memo:
\memoanswer{
（1）根据速度位移公式有 $v^{2}=2ax$

代入数据可得 $a = 2.7 \Umsq$

（2）根据向心加速度的表达式 $a=\frac{v^{2}}{R}$

可得甲、乙的向心加速度之比为

$\frac{a_{\text{甲}}}{a_{\text{乙}}}=\frac{v_{\text{甲}}^{2}}{v_{\text{乙}}^{2}}\times\frac{R_{\text{乙}}}{R_{\text{甲}}}=\frac{225}{242}$

甲、乙两物体做匀速圆周运动，则运动的时间为 $t=\frac{\pi R}{v}$

代入数据可得甲、乙运动的时间为 $t_{\text{甲}}=\frac{4\pi}{5}\Us$，$t_{\text{乙}}=\frac{9\pi}{11}\Us$

因 $t_{\text{甲}} < t_{\text{乙}}$，所以甲先出弯道。
}

\item
%% number: 14
%% paperName: 2022年辽宁高考第 14 题 12 分
%% typeId: 6
%% type: 计算题
%% chapter: 电场
%% point: 带电粒子的偏转
%% method: 无
%% score: 12
%% degree: 550
%% duplicateId: 0
%% body:
如图所示，光滑水平面 AB 和竖直面内的光滑 $\frac{1}{4}$ 圆弧导轨在 B 点平滑连接，导轨半径为 $R$。质量为 $m$ 的带正电小球将轻质弹簧压缩至 A 点后由静止释放，脱离弹簧后经过 B 点时的速度大小为 $\sqrt{gR}$，之后沿轨道 BO 运动。以 $O$ 为坐标原点建立直角坐标系 $xOy$，在 $x \geq -R$ 区域有方向与 $x$ 轴夹角为 $\theta = 45{^\circ}$ 的匀强电场，进入电场后小球受到的电场力大小为 $\sqrt{2}mg$。小球在运动过程中电荷量保持不变，重力加速度为 $g$。求：

\begin{enumerate}
	\item
	弹簧压缩至 A 点时的弹性势能；
	\item
	小球经过 O 点时的速度大小；
	\item
	小球过 O 点后运动的轨迹方程。
\end{enumerate}

\onepicture[w=5.20cm, align=right]{figs/14.png}

%% answer:
\jdanswer{
\begin{enumerate}
	\item
	$\frac{1}{2}mgR$

	\item
	$v_{O} = \sqrt{3gR}$

	\item
	$y^{2} = 6Rx$
\end{enumerate}
}
%% memo:
\memoanswer{
（1）小球从 A 到 B，根据能量守恒定律得

$E_{p}=\frac{1}{2}mv_{B}^{2}=\frac{1}{2}mgR$

（2）小球从 B 到 O，根据动能定理有

$-mgR + qE \cdot \sqrt{2}R = \frac{1}{2}mv_{O}^{2} - \frac{1}{2}mv_{B}^{2}$

解得 $v_{O} = \sqrt{3gR}$

（3）小球运动至 O 点时速度竖直向上，受电场力和重力作用，将电场力分解到 x 轴和 y 轴，则 x 轴方向有

$qE\cos45^{\circ}=ma_{x}$

竖直方向有

$qE\sin45^{\circ}-mg=ma_{y}$

解得 $a_{x}=g,\quad a_{y}=0$

说明小球从 O 点开始以后的运动为 x 轴方向做初速度为零的匀加速直线运动，y 轴方向做匀速直线运动，即做类平抛运动，则有

$x=\frac{1}{2}gt^{2},\quad y=v_{O}t$

联立解得小球过 O 点后运动的轨迹方程 $y^{2}=6Rx$
}

\item
%% number: 15
%% paperName: 2022年辽宁高考第 15 题 18 分
%% typeId: 6
%% type: 计算题
%% chapter: 电磁感应
%% point: 双杆问题
%% method: 无
%% score: 18
%% degree: 350
%% duplicateId: 0
%% body:
如图所示，两平行光滑长直金属导轨水平放置，间距为 $L$。abcd 区域有匀强磁场，磁感应强度大小为 $B$，方向竖直向上。初始时刻，磁场外的细金属杆 M 以初速度 $v_{0}$ 向右运动，磁场内的细金属杆 N 处于静止状态。两金属杆与导轨接触良好且运动过程中始终与导轨垂直。两杆的质量均为 $m$，在导轨间的电阻均为 $R$，感应电流产生的磁场及导轨的电阻忽略不计。

\begin{enumerate}
	\item
	求 M 刚进入磁场时受到的安培力 $F$ 的大小和方向；
	\item
	若两杆在磁场内未相撞且 N 出磁场时的速度为 $\frac{v_{0}}{3}$，求：① N 在磁场内运动过程中通过回路的电荷量 $q$；② 初始时刻 N 到 ab 的最小距离 $x$；
	\item
	初始时刻，若 N 到 cd 的距离与第（2）问初始时刻的相同、到 ab 的距离为 $kx$（$k > 1$），求 M 出磁场后不与 N 相撞条件下 $k$ 的取值范围。
\end{enumerate}

\onepicture[w=8.20cm, align=right]{\includesvg{figs/15}}

%% answer:
\jdanswer{
\begin{enumerate}
	\item
	$F = \frac{B^{2}L^{2}v_{0}}{2R}$，方向水平向左；

	\item
	① $q = \frac{mv_{0}}{3BL}$，② $x = \frac{2mv_{0}R}{3B^{2}L^{2}}$

	\item
	$2 \leq k < 3$
\end{enumerate}
}
%% memo:
\memoanswer{
（1）细金属杆 M 以初速度 $v_{0}$ 向右刚进入磁场时，产生的动生电动势为 $E = BLv_{0}$

电流方向为 $a \to b$，电流的大小为 $I = E/2R$

则所受的安培力大小为 $F = BIL = \frac{B^{2}L^{2}v_{0}}{2R}$

安培力的方向由左手定则可知水平向左；

（2）①金属杆 N 在磁场内运动过程中，由动量定理有

$B\overline{I}L\cdot\Delta t=m\cdot\frac{v_{0}}{3}-0$

且 $q = \overline{I} \cdot \Delta t$

联立解得通过回路的电荷量为 $q = \frac{mv_{0}}{3BL}$

②设两杆在磁场中相对靠近的位移为 $\Delta x$，有 $\overline{I} = \frac{\overline{E}}{2R}$，$\overline{E} = \frac{BL \cdot \Delta x}{\Delta t}$

整理可得 $q = \frac{BL \cdot \Delta x}{2R}$

联立可得 $\Delta x = \frac{2mv_{0}R}{3B^{2}L^{2}}$

若两杆在磁场内刚好相撞，N 到 ab 的最小距离为 $x = \Delta x = \frac{2mv_{0}R}{3B^{2}L^{2}}$

（3）两杆出磁场后在平行光滑长直金属导轨上运动，若 N 到 cd 的距离与第（2）问初始时刻的相同、到 ab 的距离为 $kx$（$k > 1$），则 N 到 ab 边的速度大小恒为 $\frac{v_{0}}{3}$，根据动量守恒定律可知 $mv_{0} = mv_{1} + m \cdot \frac{v_{0}}{3}$

解得 N 出磁场时，M 的速度大小为 $v_{1}=\frac{2v_{0}}{3}$

由题意可知，此时 M 到 cd 边的距离为 $s = (k - 1)x$

若要保证 M 出磁场后不与 N 相撞，则有两种临界情况：

① M 减速到 $\frac{v_{0}}{3}$ 时出磁场，速度刚好等于 N 的速度，一定不与 N 相撞，对 M 根据动量定理有

$B\overline{I}_{1}L\cdot\Delta t_{1}=m\cdot\frac{2}{3}v_{0}-m\cdot\frac{v_{0}}{3}$

$q_{1}=\overline{I}_{1}\cdot\Delta t_{1}=\frac{BL\cdot(k-1)x}{2R}$

联立解得 $k=2$

② M 运动到 cd 边时，恰好减速到零，则对 M 由动量定理有

$B\overline{I}_{2}L\cdot\Delta t_{2}=m\cdot\frac{v_{0}}{3}-0$

同理解得 $k=3$

综上所述，M 出磁场后不与 N 相撞条件下 $k$ 的取值范围为 $2 \leq k < 3$。
}

\end{enumerate}

\end{document}
